\documentclass[aspectratio=169,notheorems]{corvinusmetropolis}
\newtheorem{theorem}{{Tétel.}}
\newtheorem{proposition}[theorem]{{Állítás.}}
\newtheorem{lemma}[theorem]{{Lemma.}}
\newtheorem{FA}[theorem]{{Az algebra alaptétele.}}
\newtheorem{corollary}[theorem]{{Következmény.}}
\theoremstyle{definition}
\newtheorem{definition}[theorem]{{Definíció.}}
\newtheorem{defprop}[theorem]{{Definíció-Állítás.}}
\newtheorem{notes}[theorem]{{Megjegyzés.}}

\email{magyarkuti@uni-corvinus.hu}
\author[Magyarkuti]{Gyula~Magyarkuti\,\orcidlink{0000-0003-4444-0481}}
\institute[Corvinus Egyetem]%
{%\inst{1}
  Budapesti Corvinus Egyetem\\
  Matematika Tanszék\\
  Sóház, 208/B

  \begin{minipage}[b]{0.8\linewidth}
  email: \putemail\\
  letölthető: \url{https://magyarkuti.github.io/algebra-1}
  \end{minipage}
  \hfill{\adjustbox{valign=B}{\qrcode[hyperlink]{https://magyarkuti.github.io/algebra-1}}}%% iff qrcode needed, otherwise comment this line!
}
\date[Tavasz]{\ontoday}%% overwrite later, if needed.

\titlegraphic{\hfill{\includegraphics[width=.2\linewidth]{\myinstlogo}}}

\renewcommand{\emph}{\alert}
\renewcommand{\mathbf}{\mathbb}
\newcommand{\lin}{\operatorname{lin}}
\newcommand{\gen}{\operatorname{gen}}
\newcommand{\im}{\operatorname{Im}}
\newcommand{\id}{\operatorname{id}}
\newcommand{\sgn}{\operatorname{sgn}}
\newcommand{\tr}{\operatorname{tr}}
\newcommand{\codim}{\operatorname{codim}}
%newcommand{\deg}{\operatorname{deg}}
\newcommand{\ip}[2]{\langle#1,#2\rangle}
\newcommand{\uj}{\text{új}}
\newcommand{\rgi}{\text{régi}}
\newcommand{\Star}[1]{#1\ensuremath{^*}\kern-\scriptspace}
\newcommand{\CStar}{\Star{\ensuremath{\mathrm{C}}}}
\endinput


\title[Algebra 1.]{Vektorterek bevezetése}
\subtitle{Altér és generátorrendszer}
\date[7. hét]{2026. november 2-án}

\begin{document}
\openingframes{Tartalom}
\section{A vektortér fogalma}
\frame{
\begin{definition}[Vektortér]\index{vektortér}
	Legyen adva egy $\mathbb{F}$ test, és egy $V$ halmaz.
	Tegyük fel, hogy adott egy $+:V\times V\to V$ két változós művelet (ezt összeadásnak nevezzük)
	és adott egy $\cdot:\mathbb{F}\times V\to V$ szintén kétváltozós művelet (ezt skalárral való szorzásnak, vagy számmal való szorzásnak mondjuk).

    A $V$-t az $\mathbb{F}$ test feletti \alert{\emph{vektortérnek}} nevezzük,
	ha $\left( V,+ \right)$ egy Abel-csoport,
	és a számmal valós szorzás műveletre teljesülnek az alábbi axiómák:
	\begin{enumerate}
		\item Minden $\alpha\in\mathbb{F}$ és minden $u,v\in V$ mellett
		      \begin{math}
			      \alpha\cdot\left( u+v \right)=\alpha\cdot u+\alpha\cdot v,
		      \end{math}
		\item Minden $\alpha,\beta\in\mathbb{F}$ és minden $u\in V$ mellett
		      \begin{math}
			      \left( \alpha+\beta \right)\cdot u=\alpha\cdot u+\beta\cdot u,
		      \end{math}
		\item Minden $\alpha,\beta\in\mathbb{F}$ és minden $u\in V$ mellett
		      \begin{math}
			      \alpha\cdot\left( \beta\cdot u \right)=\left( \alpha\beta \right)\cdot u,
		      \end{math}
		\item Minden $u\in V$ esetén $1\cdot u=u$.\qedhere
	\end{enumerate}
\end{definition}
}
\frame{
Nagyon hasonlóan ahhoz, ahogyan testben is meggondoltuk igazak a következő számolási szabályok.
Minden $\alpha\in\mathbb{F}$ mellett
\begin{enumerate}
	\item $\alpha\cdot 0=0$,
	\item $0\cdot v=0$,
	\item $\left( -1 \right)\cdot v=-v$,
	\item $\alpha\cdot v=0$ esetén $\alpha=0$ vagy $v=0$.
\end{enumerate}
A számmal való szorzás $\cdot:\mathbb{F}\times V\to V$ művelet eredményének szokásos rövidítése,
hogy a kissé körülményes $\alpha\cdot v$ helyett csak $\alpha v$-t írunk.
Az is előfordul,
különösen amikor egy konkrét vektortér konkrét műveletéről van szó,
hogy az $\alpha v$ és a $v\alpha$ jelölést is
ugyanarra az $\alpha\cdot v\in \mathbb{F}$ elemre használjuk.
}
\frame{
Alapvető példa vektortérre a mátrixok tere.
Az $m\times n$ méretű mátrixok $\mathbb{F}^{m\times n}$ halmaza a mátrix összeadással és a számmal való szorzással
vektorteret alkot az $\mathbb{F}$ test felett.
Ha $n=1$, akkor kapjuk a az oszlopvektorok $\mathbb{F}^m$ terét, ami így szintén egy $\mathbb{F}$ feletti vektortér.
Speciális esetként $\mathbb{R}^m$ egy $\mathbb{R}$ feletti vektortér,
$\mathbb{C}^m$ egy $\mathbb{C}$ feletti vektortér.

Fontos példa még, egy adott $X$ halmazból az $\mathbb{F}$ testbe képező összes függvények halmaza
a függvények közt szokásos összeadás művelettel, és számmal való szorzással.
Ezt a teret $\mathbb{F}^X$ módon szokás jelölni, és minden nehézség nélkül ellenőrizhető,
hogy $\mathbb{F}^X$ egy $\mathbb{F}$ feletti vektortér.
Hasonlóan látszik, hogy például az összes valós--valós folytonos függvények is egy $\mathbb{R}$
feletti vektorteret alkotnak, vagy ha ezek közül csak a differenciálható függvényekre szorítkozunk,
akkor ezen függvénytér is vektortér az $\mathbb{R}$ test felett.
}
\frame{
Rögzítenünk kell magunkban, hogy a vektortér definíciójában a test is fontos szerepet játszik.
Más test felett ugyanaz az Abel-csoport már egy másik vektortérhez tartozik.
Világos például, hogy $\mathbb{R}$ az $\mathbb{R}$ test felett vektortér a szokásos műveletekkel,
de látjuk majd, hogy egészen más tulajdonságai vannak annak a vektortérnek,
amit akkor kapunk,
ha az $\mathbb{R}$ valós számok Abel-csoportját,
mint a $\mathbb{Q}$ racionális test feletti vektortérnek tekintjük.

Játékos példaként gondoljuk meg,
hogy a pozitív valós számok egy az $\mathbb{R}$ feletti vektorteret alkot a következő fura műveletekkel:
Tetszőleges $a,b$ pozitív valós szám mellett
\(
a\#b=a\cdot b,
\)
majd tetszőleges $\alpha$ valós szám és $a$ pozitív valós számra legyen
\(
\alpha\ast a=a^\alpha.
\)
\footnote{Itt $a\ast\alpha$ mit jelent?}
}
\frame{
A vektorteret sokszor csak az additív művelet alaphalmazával jelöljük.

Kicsit pontosabb ha az alaphalmaz mellett a testet is konkrétan specifikáljuk,
de sokszor feltesszük, hogy a szövegkörnyezetből nyilvánvaló, hogy mely testre gondolunk.

Hasonlóan pontosabb lenne a két műveletet is mindig kijelölni, mikor egy vektortérre hivatkozunk,
de ha világos, hogy mi a vektortérbeli elemek közt az additív művelet, és hogy mit jelent egy test elemeivel szorozni,
akkor elhagyjuk a műveletek kijelölését.

A legpontosabb, -- de persze a legkörülményesebb -- jelölés az lenne, hogy pl.
,,tekintsünk egy $\left( V,+,\cdot \right)$ vektorteret az $\mathbb{F}$ test fölött.''
Ehelyett sokszor csak azt mondjuk, hogy
,,legyen $V$ egy vektortér''.
Ilyenkor a szövegkörnyezetből világosnak kell lennie, hogy mi a test, mit jelent a számmal való szorzás, és mi az összeadás a $V$ halmazon.
}

\section{Vektortér alterei}
\frame{
\begin{definition}[altér]\index{vektortér altere}\index{altér}
	Legyen $\left( V,+,\cdot \right)$ egy vektortér az $\mathbb{F}$ test felett.
	Egy $M\subseteq V$ részhalmaz a $V$ vektortér \emph{altere},
	ha $M$ maga is vektorteret alkot a $V$-ben definiált additív művelettel,
	és a $V$-ben definiált számmal való szorzással.
\end{definition}
\begin{proposition}
	Legyen $V$ egy vektortér, és $M\subseteq V$ egy részhalmaza.
	Az $M$ pontosan akkor altér,
	ha
	\begin{enumerate}
		\item $0\in M$,
		\item $u,v\in M$ esetén $u+v\in M$,
		\item $u\in M$ és $\alpha\in\mathbb{F}$ esetén $\alpha u\in M$.\qedhere
	\end{enumerate}
\end{proposition}
Tetszőleges $V$ vektortérre a $\left\{ 0 \right\}$ és maga $V$ mindig alterek,
ezeket \alert{\emph{triviális altereknek}} is szokás mondani.\index{triviális altér}
}
\frame{
\begin{defprop}[generált altér]\index{generált altér}
	Egy vektortérben, akárhány altér közös része altér.
	Emiatt értelmes a következő definíció.
	Ha $H\subseteq V$ egy részhalmaza a $V$ vektortérnek,
	akkor jelölje
	$\gen H$ a $H$ halmazt tartalmazó összes alterek metszetét.

	Ezt az alteret nevezzük a $H$ halmaz által \emph{generált altérnek.}
\end{defprop}
\metroset{block=fill}
\begin{proposition}
	Legyen $V$ egy vektortér.
	Ekkor
	\begin{enumerate}
		\item Minden $H\subseteq V$ halmazra $H\subseteq \gen H$,
		\item Ha $H\subseteq K\subseteq V$, akkor $\gen H\subseteq \gen K$,
		\item Minden $H\subseteq V$ mellett $\gen\left( \gen H \right)=\gen H$.
	\end{enumerate}
	A $\gen H$ a $H$ halmazt tartalmazó alterek közt a legszűkebb.

	Így $H\subseteq V$ pontosan akkor altér, ha $\gen H=H$.
\end{proposition}
\metroset{block=transparent}
}
\frame{
\begin{definition}[lineáris kombináció, lineáris burok]\index{lineáris kombináció}\index{lineáris burok}
	Ha adott a $V$ vektortérben véges sok $v_1,\dots,v_r$ vektor,
	akkor a vektortér minden
\metroset{block=fill}
\begin{block}{}
	\[
		\alpha_1v_1+\dots+\alpha_r v_r
	\]
\end{block}
\metroset{block=transparent}
	alakú vektorát a $v_1,\dots,v_r$ vektorok egy \emph{lineáris kombinációjának} mondjuk.

	Legyen $H\subseteq V$ egy tetszőleges halmaz.
	Ekkor $\lin H$ jelöli $H$ összes véges részhalmazának összes lineáris kombinációinak halmazát,
	azaz
\metroset{block=fill}
\begin{block}{}
	\[
		\lin H=
		\left\{ \sum_{j=1}^n\alpha_jv_j:n\in\mathbb{N},v_1,\dots,v_n\in H,\alpha_1,\dots,\alpha_n\in\mathbb{F} \right\}
	\]
\end{block}
\metroset{block=transparent}
	Definíció szerint nulla darab vektor lineáris kombinációja a vektortér zérus eleme,
	tehát $\lin \emptyset=\left\{ 0 \right\}.$
	A $\lin H$ halmazt nevezzük a $H$ halmaz \emph{lineáris burkának}.
\end{definition}
}
\frame{
\begin{proposition}
	Egy $V$ vektortér minden $H\subseteq V$ részhalmazának lineáris burka,
	a $H$ halmazt tartalmazó legszűkebb altér, azaz
	\[
		\lin H=\gen H.\qedhere
	\]
\end{proposition}

%\begin{proof}
%	Világos, hogy $H\subseteq \lin H$, világos hogy $\lin H$ egy altér,
%	és az is nyilvánvaló, hogy ha $H\subseteq M$ egy tetszőleges altér,
%	akkor $\lin H\subseteq M$.
%
%	A $\lin H$ egy a $H$-t tartalmazó altér,
%	és $\gen H$ az összes ilyen alterek metszete, ezért
%	$\gen H\subseteq\lin H$.
%	Másoldalról, a $\lin H$ altér részhalmaza minden $H$-t tartalmazó altérnek,
%	így azok közös részének is,
%	ergo
%	$\lin H\subseteq \gen H$.
%\end{proof}
Ha $H$ egy véges halmaz, akkor $\lin H$ kicsit egyszerűbben írható.
Mint arról már a definícióban is szó volt $\lin \emptyset=\left\{ 0 \right\}$.

Ha $H$ egy elemű, akkor $\lin (\left\{ v_1 \right\})=\left\{ \alpha v_1:\alpha\in\mathbb{F} \right\}$.

Ha $H=\left\{ v_1,v_2 \right\}$ két elemű, akkor
$\lin({v_1,v_2})=\left\{ \alpha_1v_1+\alpha_2v_2:\alpha_1,\alpha_2\in\mathbb{F} \right\}$.

Hasonlóan, ha $H=\left\{ v_1,\dots,v_r \right\}$ halmaz $r$ elemből áll akkor elegendő csak az $r$ elemből álló lineáris kombinációkat képezni, azaz
\[
	\lin \left( \left\{ v_1,\dots,v_r \right\} \right)
	=
	\left\{ \sum_{j=1}^r\alpha_jv_j:\alpha_1,\dots,\alpha_r\in\mathbb{F} \right\}.
\]
}
\frame{
\begin{definition}[generátorrendszer, végesen generált vektortér]\index{generátorrendszer}\index{végesen generált vektortér}
	Egy vektortér egy $H$ részhalmazáról azt mondjuk, hogy \emph{generálja a vektorteret}
	vagy, hogy $H$ egy \emph{generátorrendszere} $V$-nek, ha
	\[
		\lin H=V.
	\]
	A $V$ vektorteret \emph{végesen generáltnak} mondunk, ha létezik véges generátorrendszere.
\end{definition}
}
\frame{
\begin{lemma}[generátorrendszer csere]\label{le:gencsere}
	Legyen $\left\{ x_1,\dots,x_m \right\}$ egy generátorrendszere valamely vektortérnek,
	és tegyük fel, hogy valamely $y$ vektorra
	\[
		y=\sum_{j=1}^m\eta_jx_j,
	\]
	ahol $\eta_k\neq 0$ valamely $1\leq k\leq m$ mellett.
	Ekkor $y$ becserélhető a $k$-adik helyen a generátorrendszerbe,
	úgy hogy az generátorrendszer maradjon, azaz a
	\[
		\left\{ x_1,\dots,x_{k-1},y,x_{k+1},\dots,x_m \right\}
	\]
	vektorrendszer is generátorrendszer.
\end{lemma}
}
%\begin{proof}
%	Fejezzük ki $x_k$-t az $y$-ra felírt formulából:
%	\(\displaystyle
%	x_k=\frac{1}{\eta_k}y+\sum_{\substack{j=1\\j\neq k}}^m\frac{-1}{\eta_k}\eta_jx_j.
%	\)
%	Ha $a$ eredetileg
%	\(\displaystyle
%	a=\sum_{j=1}^m\alpha_jx_j
%	\)
%	alakú, akkor $x_k$ helyére betéve, a fent kifejezett formulát és bevezetve a
%	$\delta=\frac{\alpha_k}{\eta_k}$ jelölést, azt kapjuk hogy:
%	\begin{multline*}
%		a=\alpha_kx_k+\sum_{\substack{j=1\\j\neq k}}^m\alpha_jx_j=
%		\\
%		=
%		\alpha_k
%		\left(
%		\frac{1}{\eta_k}y+\sum_{\substack{j=1\\j\neq k}}^m\frac{-1}{\eta_k}\eta_jx_j
%		\right)
%		+\sum_{\substack{j=1\\j\neq k}}^m\alpha_jx_j
%		=
%		\frac{\alpha_k}{\eta_k}y+
%		\sum_{\substack{j=1\\j\neq k}}^m\left( \alpha_j-\frac{\alpha_k}{\eta_k}\eta_j \right)x_j=
%		\\
%		=\delta y+
%		\sum_{\substack{j=1\\j\neq k}}^m\left( \alpha_j-\delta\eta_j \right)x_j.
%	\end{multline*}
%	Azt kaptuk tehát, hogy ha egy vektor kifejezhető az eredeti vektorrendszerből az
%	\[
%		\left( \alpha_1,\dots,\alpha_m \right)
%	\]
%	együtthatókkal, akkor ugyanez a vektor a módosított vektorrendszerből is kifejezhető,
%	méghozzá az
%	\[
%		\left(
%		\underbrace{\alpha_1-\delta\eta_1}_{1.},
%		\underbrace{\alpha_2-\delta\eta_2}_{2.},
%		\dots,
%		\underbrace{\alpha_{k-1}-\delta\eta_{k-1}}_{k-1.},
%		\underbrace{\delta}_{k.},
%		\underbrace{\alpha_{k+1}-\delta\eta_{k+1}}_{k+1.},\dots,
%		\underbrace{\alpha_m-\delta\eta_m}_{m.}
%		\right)
%	\]
%	együtthatókkal.
%\end{proof}
\frame{
	Azt kaptuk tehát, hogy ha egy vektor kifejezhető az eredeti vektorrendszerből az
	\[
		\left( \alpha_1,\dots,\alpha_m \right)
	\]
	együtthatókkal, akkor ugyanez a vektor a módosított vektorrendszerből is kifejezhető,
	méghozzá az
	\[
		\left(
		\underbrace{\alpha_1-\delta\eta_1}_{1.},
		\underbrace{\alpha_2-\delta\eta_2}_{2.},
		\dots,
		\underbrace{\alpha_{k-1}-\delta\eta_{k-1}}_{k-1.},
		\underbrace{\delta}_{k.},
		\underbrace{\alpha_{k+1}-\delta\eta_{k+1}}_{k+1.},\dots,
		\underbrace{\alpha_m-\delta\eta_m}_{m.}
		\right)
	\]
}
\section{Elimináció}
\frame{
Szokásos jelölés és a kívánatos szemlélet kialakításában is fontos szerepet játszik a következő táblázat.
Ha $\left\{ x_1,\dots,x_m \right\}$ egy generátorrendszer az azt jelenti, hogy minden $a\in V$ vektor
előáll mint az $x_1,\dots,x_m$ vektorok valamilyen együtthatókkal vett lineáris kombinációja.
Ha tehát $a=\alpha_1x_1+\dots+\alpha_mx_m$, akkor azt a következőképpen fejezzük ki.
\[
	\begin{array}{c|c}
		       & a        \\
		\hline
		x_1    & \alpha_1 \\
		\vdots & \vdots   \\
		x_k    & \alpha_k \\
		\vdots & \vdots   \\
		x_m    & \alpha_m \\
		\hline
	\end{array}
\]
Tekinthető ez egy $m\times 1$ típusú bekeretezett mátrixnak, a hol a sorok címkéi a generátorrendszer elemei,
az egyetlen oszlop címkéje pedig az a vektor, amelynek az előállításáról van szó.
}
\frame{
A generátorrendszer csere lemma arról is szólt, hogy ha $y$-t a $k$-adik helyen cseréljük a generátorrendszerbe,
akkor a fenti táblázat hogyan változik.
Azt kaptuk, hogy a
\begin{eqnarray}\label{alg:G-J}
	\begin{array}{c|cc}
		       & y              & a                       \\
		\hline
		x_1    & \eta_1         & \alpha_1                \\
		\vdots & \vdots         & \vdots                  \\
		x_k    & \boxed{\eta_k} & \alpha_k                \\
		\vdots & \vdots         & \vdots                  \\
		x_m    & \eta_m         & \alpha_m                \\
		\hline
		       & \delta         & \frac{\alpha_k}{\eta_k}
	\end{array}
	&\implies&
	\begin{array}{c|cc}
		       & y      & a                     \\
		\hline
		x_1    & 0      & \alpha_1-\eta_1\delta \\
		\vdots & \vdots & \vdots                \\
		y      & 1      & \delta                \\
		\vdots & \vdots & \vdots                \\
		x_m    & 0      & \alpha_m-\eta_m\delta \\
		\hline
		       &        &
	\end{array}
\end{eqnarray}
transzformációt kell végrehajtani.}
\frame{
Persze ugyanezt azt egyetlen $a$ vektor helyett kiszámolhatjuk például az $a,b,c$ vektorokra is.
Továbbra is az $y$-t cseréljük a generátorrendszerbe, így a
\begin{eqnarray*}
	\begin{array}{c|cccc}
		       & y              & a                       & b                      & c                       \\
		\hline
		x_1    & \eta_1         & \alpha_1                & \beta_1                & \gamma_1                \\
		\vdots & \vdots         & \vdots                  & \vdots                 & \vdots                  \\
		x_k    & \boxed{\eta_k} & \alpha_k                & \beta_k                & \gamma_k                \\
		\vdots & \vdots         & \vdots                  & \vdots                 & \vdots                  \\
		x_m    & \eta_m         & \alpha_m                & \beta_m                & \gamma_m                \\
		\hline
		       & \delta         & \frac{\alpha_k}{\eta_k} & \frac{\beta_k}{\eta_k} & \frac{\gamma_k}{\eta_k} \\
	\end{array}
	&\implies&
	\begin{array}{c|cccc}
		       & y      & a                       & b                      & c                       \\
		\hline
		x_1    & 0      & \alpha_1-\eta_1\delta_a & \beta_1-\eta_1\delta_b & \gamma_1-\eta_1\delta_c \\
		\vdots & \vdots & \vdots                  & \vdots                 & \vdots                  \\
		y      & 1      & \delta_a                & \delta_b               & \delta_c                \\
		\vdots & \vdots & \vdots                  & \vdots                 & \vdots                  \\
		x_m    & 0      & \alpha_m-\eta_m\delta_a & \beta_m-\eta_m\delta_b & \gamma_m-\eta_m\delta_c \\
		\hline
		       &        &
	\end{array}
\end{eqnarray*}
a transzformációt hajtjuk végre.
}
\frame{
Most arra figyeljünk,
hogy végül is a keretben lévő $m\times 4$-es mátrixszal sorműveleteket hajtottunk végre:
\begin{itemize}
	\item
	      Eldöntöttük, hogy az $y,a,b,c$ vektorok közül az $y$-t cseréljük be a generátorrendszerbe,
	      mégpedig a $k$-adik helyen.
	      Ezt jeleztük az $y$ oszlopa $k$-adik helyen lévő elemének bekeretezésével.
	      A keretben csak nem zéró szám lehet.
	\item
	      Első lépésként a keretezett számmal osztottuk a $k$-adik sort. Ez az új táblázat $k$-adik sora.
	      Pusztán segítségképpen ugyanezt a sort mint egy számolási segédsort az első táblázat alá másoltuk.
	\item
	      Az új táblázat első sora az eredeti első sornak és a segédsor $\eta_1$-szeresének különbsége.
	      A második sor az eredeti második sornak és a segédsor $\eta_2$-szeresének különbsége.
	      Hasonlóan, $m\neq k$-ra az $m$-edik sor az eredeti $m$-edik sornak és a segédsor $\eta_m$-szeresének különbsége.
\end{itemize}
}
\frame{
Látjuk tehát, hogy összesen két fajta sorművelettel alakítottuk a kiindulási mátrixot:
\begin{enumerate}
	\item Egy sort szoroztunk egy nem zérus számmal,
	\item Egy sorhoz hozzáadtuk egy másik sor számszorosát.
\end{enumerate}
Világos, hogy ez ugyanaz, mintha a feladat az lett volna,
hogy a fenti két sorművelet használva az $y$ oszlopában minden számot el kell tüntetni,
a $k$-adikat pedig $1$-re kell beállítani.
Ez a \emph{Gauss\,--\,Jordan-elimináció}.\index{Gauss\,--\,Jordan-elimináció}
}
%\closingframes{Köszönöm a figyelmet, a türelemet, a megértést, az empátiát!}
\frame{}
\end{document}

